\chapter{Абетка Морзе}
% full version: chapters/04_code.tex

В абетці Морзе два знаки --- крапка й тире --- і паузи між ними. У нейрона алфавіт ще бідніший: знак один, коротке клацання. Отже, усе повідомлення записано в паузах. Якою ж мовою розмовляють нейрони?

\section*{Скільки можна сказати клацаннями}

Почнімо як інженер зв'язку: скільки інформації взагалі вміщується в такому каналі? Якщо отримувач розрізняє час приходу клацання з точністю до мілісекунди, одне клацання може нести близько восьми бітів. Якщо ж отримувач лише рахує, скільки клацань прийшло за секунду, --- з того самого потоку він видобуде в десятки разів менше. Майже вся ємність каналу сидить у \emph{часі}.

\pencilfig{4}{\pencilart{4}}{Згори абетка Морзе, знизу нейрон: знак один, клацання; усе повідомлення записано в паузах.}

Але канал шумний. Сам нейрон, як не дивно, точний: якщо багато разів подати на нього той самий змінний струм, він клацатиме в ті самі моменти з точністю до мілісекунди. Шумлять з'єднання: понад половина клацань, що добігли до кінця проводу, просто губиться --- речовина викидається не щоразу. У живій корі потік клацань виглядає майже випадковим. Чи це шум --- чи повідомлення, яке ми поки не вміємо прочитати? Добре стиснуте повідомлення для стороннього не відрізнити від випадкового, тож питання відкрите.

\section*{Повільно, але надійно}

Найпростіший код --- рахувати клацання: що частіше, то більше число. Цьому кодові не страшні ні втрати, ні тремтіння. Але він повільний. Щоб дізнатися частоту з точністю до десяти відсотків, потрібна сотня клацань, а це секунди. Тим часом ви впізнаєте тварину на картинці, показаній на мить, приблизно за півтори десятої секунди, і сигнал при цьому проходить близько десятка щаблів обробки. На кожному щаблі кожен нейрон устигає клацнути щонайбільше раз.

Вихід --- узяти багато нейронів одразу: не один слухач рахує клацання хвилину, а сотня слухачів --- мить. Тільки й тут каверза: шуми сусідніх нейронів трохи схожі, як чутки в натовпі. Усереднення за кількома десятками сусідів допомагає, а далі майже перестає.

\pencilfig{122}{%
% error of the group-average rate relative to one neuron, sqrt(c + (1-c)/N): c = 0.1 (neighbours' noise
% is slightly alike; full edition, chapter 4) and c = 0 (independent noise). x = 0.8 log2 N, y = 2.2 * error
\draw[pen] (0,0) -- (5.9,0); \draw[pen] (0,0) -- (0,2.45);
\foreach \x/\n in {0/1, 2.66/10, 5.31/100}{\draw[pcGray, line width=0.5pt] (\x,0) -- (\x,-0.08); \node[lbl, anchor=north] at (\x,-0.08) {\n};}
\node[lbl, anchor=north] at (2.95,-0.38) {скільки нейронів усереднюємо};
\node[lbl, anchor=south, rotate=90] at (-0.1,1.2) {похибка};
\draw[pcGray, densely dashed, line width=0.5pt] (0,0.70) -- (5.6,0.70);
\draw[penlite=pcBlue] plot[smooth] coordinates {(0.00,2.20) (0.47,1.80) (0.80,1.56) (1.27,1.27) (1.60,1.10) (2.07,0.90) (2.40,0.78) (2.87,0.64) (3.20,0.55) (3.67,0.45) (4.00,0.39) (4.47,0.32) (4.80,0.28) (5.27,0.22) (5.60,0.19)};
\draw[penlite=pcRed] plot[smooth] coordinates {(0.00,2.20) (0.47,1.84) (0.80,1.63) (1.27,1.39) (1.60,1.25) (2.07,1.10) (2.40,1.01) (2.87,0.92) (3.20,0.87) (3.67,0.82) (4.00,0.79) (4.47,0.76) (4.80,0.74) (5.27,0.73) (5.60,0.72)};
\node[lbl, text=pcRed, anchor=west, align=left] at (5.7,0.74) {схожі шуми,\\як у корі};
\node[lbl, text=pcBlue, anchor=west] at (5.7,0.19) {незалежні шуми};
\node[lbl, anchor=west] at (0.9,2.15) {один нейрон};
\node[lbl, anchor=north west] at (0.05,0.66) {утричі менша};
}{Похибка частоти, усередненої за групою нейронів. Якби шуми були незалежні, вона падала б і далі. Схожі шуми сусідів зупиняють її приблизно на третині, і цієї межі сягають уже на кількох десятках нейронів.}

\section*{Швидко: хто перший}

Другий вихід --- записувати число в часі. Сильний сигнал змусить нейрон клацнути рано, слабкий --- пізно. Одне-єдине клацання вже несе число. Але від чого відлічувати час, якщо немає годинника? Можна порівнювати нейрони між собою: хто клацнув першим, хто другим. Порядок приходу клацань десяти нейронів несе понад двадцять бітів. А ще позначкою початку слугує спільний залп --- як довга пауза між словами в абетці Морзе.

\pencilfig{121}{%
% latency code: one click per neuron; the stronger the input, the earlier the click
\foreach \y/\t/\l/\k in {1.5/1.1/{сильний вхід}/1, 0.95/2.4/{середній}/2, 0.4/4.2/{слабкий}/3}{
  \draw[pen] (0.4,\y) -- (5.1,\y);
  \draw[pcBlue, line width=0.9pt] (\t,\y) -- (\t,\y+0.42);
  \node[lbl, anchor=east] at (0.3,\y+0.1) {\l};
  \draw[dotpen=pcAmber, wash=pcAmber] (5.6,\y+0.16) circle (0.17);
  \node[font=\scriptsize, text=pcAmber!60!black] at (5.6,\y+0.16) {\k};}
\draw[pcGray, densely dashed, line width=0.5pt] (0.55,0.25) -- (0.55,2.1);
\node[lbl, anchor=south] at (0.55,2.1) {подразник};
\node[lbl, text=pcAmber!70!black, anchor=south] at (5.6,2.1) {порядок};
\draw[arr] (0.7,0.05) -- (4.9,0.05); \node[lbl, anchor=north] at (2.8,0.02) {час};
}{Сильний вхід --- раннє клацання, слабкий --- пізнє. Порядок клацань показує, який вхід сильніший, навіть якщо момент початку невідомий.}

\section*{Код обирає слухач}

Та сама послідовність клацань несе і частоту, і часи, і порядок. Що з цього буде прочитано, вирішує отримувач. Нейрон із «широкою діркою» у відрі повільно накопичує воду й чує лише частоту. Нейрон зі швидким забуванням чує лише збіги. А гальма, як ми бачили, вміють перемикати нейрон з одного режиму в інший.

Навіть з'єднання фільтрують сигнал. Одні швидко втомлюються і тому повідомляють не частоту, а її \emph{зміну}. Інші, навпаки, краще пропускають пачки з кількох клацань поспіль: у нейрона з'являється «тире». Гальмівні нейрони тим часом розставляють розділові знаки --- ріжуть потік на вікна й приглушують його, коли він надто гучний.

\section*{Ощадлива мова}

І останнє: ця мова дорога. Кожне клацання коштує енергії, а весь мозок споживає близько двадцяти ватів. Якщо поділити ці вати на всі нейрони, вийде в середньому порядку одного клацання на секунду на нейрон. Більшість нейронів більшу частину часу мовчать. Код мозку мусить бути скупим.

В абетці Морзе найчастіша літера англійської мови, E, --- це одна крапка: часте кодують коротко. У нейронів схоже: те, що очікуване, коштує мало клацань. Сталий подразник поступово перестає викликати відповідь, датчики повідомляють про зміни, а дофамін --- лише про сюрпризи. Мозок витрачає свої клацання на новини.

\fullversion{4}
